\documentclass[11pt]{report}

% ---------------------------------------------------------------------------
% AFFINE JUMP INTENSITY FOR THE TODOROV TWO-FACTOR VOLATILITY MODEL
% Working notes: deriving the affine-intensity extension and the moments the
% GMM estimator needs.
%
% Compile with pdflatex (twice, for TOC and references).
% ---------------------------------------------------------------------------

\usepackage[margin=1.1in]{geometry}
\usepackage{amsmath, amssymb, amsthm, mathtools}
\usepackage{booktabs}
\usepackage{xcolor}
\usepackage[colorlinks=true, linkcolor=blue!50!black, urlcolor=blue!50!black]{hyperref}

\newtheorem{theorem}{Theorem}[chapter]
\newtheorem{lemma}[theorem]{Lemma}
\newtheorem{proposition}[theorem]{Proposition}
\newtheorem{corollary}[theorem]{Corollary}
\theoremstyle{definition}
\newtheorem{definition}[theorem]{Definition}
\newtheorem{assumption}[theorem]{Assumption}
\newtheorem{example}[theorem]{Example}
\theoremstyle{remark}
\newtheorem{remark}[theorem]{Remark}
\newtheorem{impact}[theorem]{Remark}

\newcommand{\E}{\mathbb{E}}
\newcommand{\Var}{\operatorname{Var}}
\newcommand{\Cov}{\operatorname{Cov}}
\newcommand{\R}{\mathbb{R}}
\newcommand{\F}{\mathcal{F}}
\newcommand{\A}{\mathcal{A}}
\newcommand{\dd}{\,\mathrm{d}}
\newcommand{\Vc}{V^{c}}
\newcommand{\Vj}{V^{j}}
\newcommand{\TPV}{\mathrm{TPV}}
\newcommand{\JV}{\mathrm{JV}}
\newcommand{\RV}{\mathrm{RV}}
\newcommand{\QV}{\mathrm{QV}}
\newcommand{\mk}{m_{k}}
\newcommand{\mkk}{m_{k^{2}}}
\newcommand{\mh}{m_{h^{2}}}
\newcommand{\mhh}{m_{h^{4}}}
\newcommand{\mkh}{m_{kh^{2}}}

\title{\textbf{Affine Jump Intensity for the Two-Factor\\ Stochastic Volatility Model}\\[1ex]
\large From point processes to the GMM moments, the filter, and the\\
variance risk premium}
\author{Agrim}
\date{July 2026}

\begin{document}
\maketitle

\begin{abstract}
\noindent
I fit the Todorov (2010) two-factor model to Bitcoin and three diagnostics I had
held back rejected it. Realized jump variance $\JV_t$ is serially correlated
when the model says it cannot be; the cross-autocovariance between $\TPV_t$ and
$\JV_t$ runs in both directions when the model allows only one; and jump
feed-through lasts about a week when the calibrated decay kills it in hours.

All three trace to the same assumption, that jumps arrive at a \emph{constant}
rate. These notes work out the smallest change that fixes it: let the intensity
be affine in the jump-variance state, $\lambda(t)=\lambda_{0}+\lambda_{1}\Vj(t^{-})$,
so each jump makes the next one more likely. Affinity is what keeps the moments
in closed form, and that is the only reason the estimator stays tractable.

I start from point processes and compensators, show that the affine class is the
tractability boundary, identify the specification with an exponential-kernel
marked Hawkes process, and derive the population moments the GMM estimator needs
using Dynkin's formula on fixed intervals. From there: identification, the
latent-state filter, the variance risk premium, and the simulator I need to
check the algebra against before running anything on real data.
\end{abstract}

\paragraph{What this document is.} These are the notes I wrote while working the
extension out, before I implemented any of it, so they read as much like a plan as a
derivation. The later chapters propose more than I ended up building: the latent-state
filter of Chapter~\ref{ch:filter} and the regime classification in Chapter~\ref{ch:vrp}
were never written. The repository README and \texttt{writeup.pdf} record what actually
happened, and two things there went against what I predict below. The event-study
prediction, that post-event decay should match the estimated $\kappa_{j}$, failed:
observed persistence was about 9.5 days against a predicted 2.4. And the branching ratio
$\eta$ came out statistically indistinguishable from zero, so the jump clustering itself
is on much firmer ground than the structural parameter that decomposes it. The
derivations I did check, against an exact simulator, before using any of them.

\tableofcontents

% ===========================================================================
\chapter{Orientation: what broke, and what we are building}
\label{ch:orientation}

\section{The rejected specification}

The Step-A model drives price jumps and jump-variance jumps with a single
Poisson random measure $\mu(\dd t,\dd x)$ whose compensator is
$\dd t\,G(\dd x)$: \emph{time-homogeneous}. Consequences that the data
rejected:

\begin{enumerate}
\item \textbf{$\JV$ autocovariance.} Under time-homogeneity, jump energies
collected over disjoint days are independent, so
$\Cov(\JV_t,\JV_{t-k})=0$ for every lag $k\ge 1$ and every parameter value.
Empirically the lag-1 autocovariance is $\approx 12.4$ against a
$\pm 3.5$ noise band, decaying slowly through lag $\sim$10. \emph{Jumps
cluster.}
\item \textbf{Cross-autocovariance asymmetry.} The model implies
$\Cov(\TPV_t,\JV_{t+k})=0$ for $k\ge 1$: past diffusive variance cannot
predict future jump energy, because future jump arrivals are independent of
everything. Empirically that branch is clearly positive out to $k\approx 7$.
\item \textbf{Feed-through persistence.} The calibrated jump-factor decay
$|\hat\rho_j|=6.14$ (half-life 0.11 d) kills the model-implied feed-through of
jump energy into future $\TPV$ within two days; empirically it persists about
a week.
\end{enumerate}

\section{The repair in one paragraph}

Let the arrival rate of jumps be an \emph{affine function of the jump-variance
state itself}:
\[
\lambda(t) \;=\; \lambda_{0} + \lambda_{1}\,\Vj(t^{-}), \qquad
\lambda_{0}>0,\ \lambda_{1}\ge 0 .
\]
Every jump raises $\Vj$, which raises the arrival rate of \emph{further}
jumps: a self-exciting cascade. One new parameter ($\lambda_{1}$)
simultaneously produces (i) exponentially decaying $\JV$ autocovariance,
(ii) a two-sided cross-autocovariance whose forward branch exceeds its
backward branch by exactly the co-jump coupling, and (iii) an
\emph{effective} persistence rate $\kappa_{j}$ that is \emph{slower} than the
per-jump decay rate $b=-\rho_j$ --- resolving the paradox of
fast individual decay coexisting with slow observed persistence. Because the
intensity is affine, all unconditional and conditional moments remain in
closed form, so the entire GMM/filtering/VRP pipeline survives with formulas
swapped, not rebuilt.

\section{How this document is organized}

Chapter~\ref{ch:pp} builds point processes, random measures, compensators and
the martingale calculus we need --- from zero. Chapter~\ref{ch:affine}
explains what ``affine'' means, proves why it is the tractability frontier,
and identifies our specification with a marked Hawkes process.
Chapter~\ref{ch:model} states the extended model and proves the central
effective-rate theorem. Chapter~\ref{ch:moments} is the moment engine: every
population quantity the estimator needs, derived with Dynkin's formula over
fixed intervals, each ending in a boxed formula. Chapter~\ref{ch:gmm} treats
identification and the GMM design. Chapter~\ref{ch:filter} rebuilds the
latent-state projection. Chapter~\ref{ch:vrp} connects the calibrated
physical dynamics to the variance risk premium and regime classification.
Chapter~\ref{ch:impl} specifies simulation, verification, and the changes
needed to the existing estimation code. Appendices collect derivatives for the analytic
Jacobian, useful integrals, notation, and an implementation checklist.

\paragraph{How the derivations run.} Throughout, I work through \emph{Dynkin's formula evaluated over a fixed
global time interval}. I never shrink a time step to zero as a derivation device, and never reach
for the adjoint (forward/Fokker--Planck) operator.
Stationary moments follow from the fixed-interval identity
$\E\!\int_{0}^{T}\!\A f(X_{s})\dd s = \E f(X_{T})-\E f(X_{0}) = 0$;
conditional moments follow from the same identity read as a Volterra integral
equation in the interval endpoint.

% ===========================================================================
\chapter{Point processes, random measures, and compensators}
\label{ch:pp}

\section{Counting processes and the Poisson process}

\begin{definition}[Counting process]
A \emph{counting process} $N=(N_t)_{t\ge 0}$ is a right-continuous,
nondecreasing, integer-valued process with $N_0=0$ and unit jumps. $N_t$
counts events in $(0,t]$.
\end{definition}

\begin{definition}[Poisson process]
$N$ is a \emph{Poisson process with rate} $\lambda>0$ if it has independent
increments and $N_t-N_s\sim\mathrm{Poisson}(\lambda(t-s))$ for $s<t$.
\end{definition}

Two properties matter for us. First, $N_t-\lambda t$ is a martingale: the
``compensated'' process has no predictable trend. Second, independent
increments are exactly what the data rejected --- event counts (and energies)
in disjoint windows are uncorrelated. Everything in this chapter generalizes
these two facts to state-dependent rates.

\section{Marked point processes and random measures}

Our jumps carry a \emph{mark} $x$ (the jump size information): the $i$-th
event is a pair $(\tau_i, x_i)$ of a time and a mark in $\R_0=\R\setminus\{0\}$
(or a more general mark space; nothing below depends on it). The bookkeeping
device is the \emph{random counting measure}
\[
\mu(A\times B) \;=\; \#\{\,i:\ \tau_i\in A,\ x_i\in B\,\},
\qquad A\subseteq[0,\infty),\ B\subseteq\R_0 ,
\]
so that for any function $\psi$ of the mark,
\[
\int_{0}^{t}\!\!\int_{\R_0}\psi(x)\,\mu(\dd s,\dd x)
\;=\;\sum_{i:\ \tau_i\le t}\psi(x_i)
\]
is the running sum of $\psi$ over the jumps so far. Realized jump variance is
exactly such a sum with $\psi=h^{2}$; the jump-variance factor $\Vj$ is
driven by such a sum with $\psi=k$.

\section{Stochastic intensity and the compensator}

\begin{definition}[Stochastic intensity]\label{def:intensity}
Let $(\F_t)$ be the information filtration. A nonnegative,
\emph{predictable} process $\lambda(t)$ is the ($\F_t$-)\emph{intensity} of
the marked point process $\mu$ with i.i.d.\ marks $x_i\sim G$ if
\[
\nu(\dd t,\dd x)\;=\;\lambda(t)\,\dd t\,G(\dd x)
\]
\emph{compensates} $\mu$: for every bounded predictable integrand
$\psi$,
\begin{equation}\label{eq:compensation}
M_{t}\;=\;\int_{0}^{t}\!\!\int \psi(x)\,\big[\mu(\dd s,\dd x)-\nu(\dd s,\dd x)\big]
\quad\text{is an $(\F_t)$-martingale.}
\end{equation}
We write $\tilde\mu = \mu-\nu$ for the compensated measure. We normalize
$G$ to be a \emph{probability} measure, so $\lambda(t)$ is the arrival
\emph{rate} of jumps and marks are i.i.d.\ draws from $G$.
\end{definition}

Intuition: over the fixed window $[t,t+\Delta]$ the expected number of jumps,
given everything known at $t$ and along the path, is
$\E\int_{t}^{t+\Delta}\lambda(s)\dd s$; the compensator is the
``predictable budget'' of jumps, and subtracting it leaves pure surprise.
Predictability of $\lambda$ (using $\Vj(t^-)$, the left limit) is not
pedantry: the rate at which the next jump arrives must be determined by the
state \emph{just before} it arrives, otherwise a jump could raise its own
arrival probability retroactively and \eqref{eq:compensation} would fail.

\begin{remark}[Normalization choice]\label{rem:normalization}
The Step-A code used the L\'evy-measure convention, absorbing the rate into
$G$: its parameters $\mathcal{I}_{h^2}=\int h^2\dd G^{\text{old}}$ etc.\ are,
in the present convention, the \emph{products}
$\mathcal{I}_{h^2}=\lambda\,\mh$, $\mathcal{I}_{h^4}=\lambda\,\mhh$,
$\mathcal{I}_{kh^2}=\lambda\,\mkh$, where
$m_{\phi}:=\int\phi(x)\,G(\dd x)$. Keeping rate and mark distribution
separate is what lets the rate become stochastic. The dictionary between old
and new parameters is given in Table~\ref{tab:dictionary}.
\end{remark}

\section{The moment toolkit for compensated integrals}

The three formulas below are the entire ``It\^o calculus'' this project needs
for jumps. Fix a day $[t-1,t]$ and write
$M[\psi]=\int_{t-1}^{t}\!\int\psi(x)\,\tilde\mu(\dd s,\dd x)$.

\begin{lemma}[First moments]\label{lem:firstmoment}
$\displaystyle \E\int_{t-1}^{t}\!\!\int\psi(x)\,\mu(\dd s,\dd x)
= m_{\psi}\,\E\!\int_{t-1}^{t}\lambda(s)\dd s$, and $\E\,M[\psi]=0$.
\end{lemma}

\begin{lemma}[It\^o isometry for marked point processes]\label{lem:isometry}
For predictable $\psi,\phi$,
\[
\E\big(M[\psi]\,M[\phi]\big)
\;=\; m_{\psi\phi}\;\E\!\int_{t-1}^{t}\lambda(s)\,\dd s .
\]
In particular $\Var(M[\psi]) = m_{\psi^{2}}\,\E\int\lambda$.
\end{lemma}

\begin{lemma}[Orthogonality of increments]\label{lem:orthogonality}
If $Z$ is $\F_{s}$-measurable and $M[\psi]$ is a compensated integral over
$(s,t]$, then $\E[Z\,M[\psi]]=0$. In particular, compensated jump integrals
over \emph{disjoint} intervals are uncorrelated, and a compensated integral is
uncorrelated with any functional of the \emph{earlier} path.
\end{lemma}

\begin{proof}[Why these hold]
Lemma~\ref{lem:firstmoment} is the definition of compensation applied to
$\psi$ and the tower property. Lemma~\ref{lem:isometry}: jumps occur at
distinct times, so the product of two jump sums over the same window is a sum
over \emph{common} jumps of $\psi(x_i)\phi(x_i)$ plus cross terms of distinct
jumps whose compensated expectations vanish; the common-jump sum has
expectation $m_{\psi\phi}\E\int\lambda$ by Lemma~\ref{lem:firstmoment}.
Lemma~\ref{lem:orthogonality} is the martingale property plus the tower
property: $\E[Z\,M]=\E[Z\,\E(M\mid\F_{s})]=0$.
\end{proof}

These lemmas are deliberately stated with a \emph{stochastic} $\lambda$: they
are exactly where the new model's extra covariance terms will come from,
because $\E\int\lambda$ ceases to be a constant and starts to covary with the
volatility path.

% ===========================================================================
\chapter{The affine principle, and why it is the tractability frontier}
\label{ch:affine}

\section{Definition and the closure property}

\begin{definition}[Affine intensity]
The intensity is \emph{affine} in the state if
$\lambda(t)=\lambda_{0}+\lambda_{1}\Vj(t^{-})$ with $\lambda_{0}>0,\
\lambda_{1}\ge0$ (and the state dynamics keep $\Vj\ge0$, which holds here
because $\Vj$ decays exponentially between nonnegative jumps).
\end{definition}

The deep reason to insist on affinity is an algebraic closure property of the
\emph{extended generator}. For our jump factor
\[
\dd \Vj(t) \;=\; \rho_{j}\Vj(t)\,\dd t + \int_{\R_0} k(x)\,\mu(\dd t,\dd x),
\qquad b:=-\rho_{j}>0,\ k\ge 0,
\]
the generator acts on a smooth $f$ as
\begin{equation}\label{eq:generator}
\A f(v) \;=\; -b\,v\,f'(v)
\;+\;\big(\lambda_{0}+\lambda_{1}v\big)\int_{\R_0}\big[f(v+k(x))-f(v)\big]\,G(\dd x).
\end{equation}
The first term is the deterministic decay; the second says: at rate
$\lambda(v)$, replace $v$ by $v+k(x)$ with $x\sim G$.

\begin{proposition}[Polynomial closure]\label{prop:closure}
If $\lambda(v)$ is affine, then $\A$ maps polynomials of degree $n$ to
polynomials of degree $n$. If $\lambda(v)$ is a polynomial of degree
$d\ge 2$, then $\A$ maps degree $n$ to degree $n+d-1 > n$.
\end{proposition}

\begin{proof}
Apply \eqref{eq:generator} to $f(v)=v^{n}$. The drift term gives $-nbv^{n}$:
degree $n$. The jump term gives
$\lambda(v)\sum_{r=0}^{n-1}\binom{n}{r}v^{r}m_{k^{n-r}}$, a polynomial of
degree $(n-1)+\deg\lambda$. With $\deg\lambda=1$ this is degree $n$; with
$\deg\lambda=d$ it is degree $n+d-1$.
\end{proof}

\begin{impact}
Proposition~\ref{prop:closure} is what makes the rest possible. Moment
derivations reduce (Chapter~\ref{ch:moments}) to computing $\E[\A f]$ for
polynomial $f$. Under affinity, the $n$-th moment equation involves only
moments of order $\le n$: a finite, triangular linear system --- closed
form, hence an analytic GMM Jacobian, hence stable optimization. Under a
quadratic intensity, the second-moment equation would involve the third
moment, the third the fourth, ad infinitum: no closed form, forcing
simulation-based estimation (orders of magnitude slower, noisy gradients,
no analytic Jacobian). So affinity is not cosmetic: it is the difference
between rederiving a formula in an afternoon and rebuilding the estimator.
\end{impact}

\section{Equivalence with a marked Hawkes process}

\begin{proposition}[Hawkes representation]\label{prop:hawkes}
Suppose $\Vj$ starts from its stationary distribution in the infinite past.
Then
\[
\Vj(t)\;=\;\sum_{\tau_i\le t} k(x_i)\,e^{-b(t-\tau_i)},
\qquad
\lambda(t)\;=\;\lambda_{0}+\lambda_{1}\!\!\sum_{\tau_i< t} k(x_i)\,e^{-b(t-\tau_i)} .
\]
That is: the pair (jump times, intensity) is precisely a \emph{marked Hawkes
process with exponential kernel} $\lambda_{1}k(x)e^{-bu}$ and baseline
$\lambda_{0}$.
\end{proposition}

\begin{proof}
Between jumps $\Vj$ satisfies $\dd\Vj=-b\Vj\dd t$, so each jump's
contribution $k(x_i)$ decays as $e^{-b(t-\tau_i)}$; summing contributions
gives the display (the transient from the infinite past vanishes).
Substituting into $\lambda=\lambda_0+\lambda_1\Vj(t^-)$ gives the Hawkes
form.
\end{proof}

The Hawkes lens supplies the two most useful intuitions in this document.

\paragraph{Cluster (branching) representation.} A Hawkes process is
statistically identical to the following cascade: \emph{immigrant} events
arrive as a plain Poisson process with rate $\lambda_{0}$; each event (of any
generation) with mark $x$ independently spawns \emph{children} as an
inhomogeneous Poisson process with rate $\lambda_{1}k(x)e^{-bu}$, $u$ its
age. The expected number of children per event is
\begin{equation}\label{eq:eta}
\eta \;=\; \int_0^\infty \lambda_1 m_k e^{-bu}\,du
\;=\; \frac{\lambda_{1}\,\mk}{b}
\qquad\text{(the \emph{branching ratio}).}
\end{equation}
A cluster started by one immigrant has expected total size
$1+\eta+\eta^{2}+\cdots=(1-\eta)^{-1}$: finite iff $\eta<1$.

\begin{assumption}[Subcriticality]\label{ass:subcritical}
$\eta=\lambda_{1}\mk/b<1$. This is necessary and sufficient for a stationary
version of $(\Vj,\lambda)$ with finite moments of all orders used below.
\end{assumption}

\paragraph{Two decay rates.} Each individual excitation decays at rate $b$
(fast, hours). But a burst is a \emph{cluster}: children arrive while the
parent's excitation wanes, grandchildren while the children's wanes. The
aggregate relaxation of the cascade is slower than any individual decay. The
next chapter makes this exact: the observable persistence rate is
$\kappa_{j}=b(1-\eta)$, and $\eta\to1$ drags $\kappa_j\to0$ regardless of
how fast $b$ is. This is the mathematical resolution of the Step-A paradox
(fitted $|\rho_j|=6.14$ versus week-long empirical feed-through): the old
model owns one decay rate and was forced to compromise; the new model owns
two.

% ===========================================================================
\chapter{The extended model}
\label{ch:model}

\section{Full specification}

\begin{assumption}[Model $\mathcal{M}'$]\label{ass:model}
Under the physical measure $\mathbb{P}$, on a filtered probability space
carrying a Brownian motion $B_{1}$ and a marked point process $\mu$ with
i.i.d.\ marks $x\sim G$ (a probability measure) and stochastic intensity
$\lambda(t)$:
\begin{align}
\dd \Vc(t) &= \kappa_{c}\big(\bar V^{c}-\Vc(t)\big)\dd t
             + \sigma_{cv}\sqrt{\Vc(t)}\,\dd B_{1}(t),
&&\text{(CIR, unchanged)} \label{eq:cir}\\[0.5ex]
\dd \Vj(t) &= \rho_{j}\Vj(t)\,\dd t + \int_{\R_0} k(x)\,\mu(\dd t,\dd x),
&&b:=-\rho_{j}>0, \label{eq:vj}\\[0.5ex]
\lambda(t) &= \lambda_{0}+\lambda_{1}\Vj(t^{-}),
&&\lambda_{0}>0,\ \lambda_{1}\ge0, \label{eq:intensity}
\end{align}
with $B_{1}\perp\mu$, $k\ge 0$, and Assumption~\ref{ass:subcritical}. Total
diffusive variance is $\sigma^{2}(t)=\Vc(t)+\Vj(t)$; the same jumps that
feed $\Vj$ carry price-jump energy $h^{2}(x)$. The daily observables over
$[t-1,t]$ are
\[
\TPV_{t}\approx\int_{t-1}^{t}\!\big(\Vc(s)+\Vj(s)\big)\dd s,
\qquad
\JV_{t}\approx\int_{t-1}^{t}\!\!\int h^{2}(x)\,\mu(\dd s,\dd x).
\]
The translation gauge $k(x)=c\,h^{2}(x)$, $c>0$, links the two jump kernels,
so mark moments satisfy $\mk=c\,\mh$, $\mkk=c^{2}\mhh$, $\mkh=c\,\mhh$.
\end{assumption}

Setting $\lambda_{1}=0$ recovers the Step-A model exactly: the extension
\emph{nests} the old one, which is essential --- it lets the data vote on
$\lambda_{1}$, and it gives a free correctness check on every new formula
(each must reduce to the old one at $\lambda_{1}=0$).

\section{The effective persistence rate}

\begin{theorem}[Stationary mean and effective rate]\label{thm:effective}
Under Assumptions \ref{ass:subcritical}--\ref{ass:model}, the stationary
means are
\[
\bar\lambda := \E[\lambda(t)]=\frac{\lambda_{0}}{1-\eta},
\qquad
\bar v := \E[\Vj(t)]=\frac{\bar\lambda\,\mk}{b}
        =\frac{\lambda_{0}\,\mk}{b(1-\eta)} ,
\]
and the conditional mean relaxes exponentially at the \emph{effective rate}
\[
\boxed{\;\kappa_{j}\;=\;b-\lambda_{1}\mk\;=\;b\,(1-\eta)\;}
\qquad\text{i.e.}\qquad
\E\big[\Vj(t+s)\,\big|\,\F_{t}\big]
=\bar v+\big(\Vj(t)-\bar v\big)e^{-\kappa_{j}s}.
\]
Consequently $\Cov\big(\Vj(t),\Vj(t+s)\big)=\Var(\Vj)\,e^{-\kappa_{j}|s|}$.
\end{theorem}

\begin{proof}
Apply Dynkin's formula on the fixed interval $[t,t+s]$ to $f(v)=v$, for
which \eqref{eq:generator} gives
$\A f(v) = -bv+(\lambda_{0}+\lambda_{1}v)\,\mk$:
\[
\E_{t}\big[\Vj(t+s)\big]-\Vj(t)
=\E_{t}\!\int_{t}^{t+s}\!\Big[-b\,\Vj(u)+\mk\big(\lambda_{0}
+\lambda_{1}\Vj(u)\big)\Big]\dd u .
\]
Write $g(s)=\E_{t}[\Vj(t+s)]$. The identity above, read for every endpoint
$s$ of the fixed interval, is the Volterra equation
$g(s)=g(0)+\int_{0}^{s}\!\big[\lambda_{0}\mk-\kappa_{j}\,g(u)\big]\dd u$
with $\kappa_{j}=b-\lambda_{1}\mk$; its unique solution is the boxed
exponential with resting point $\lambda_{0}\mk/\kappa_{j}=\bar v$.
Stationary means: take $f(v)=v$ in the stationarity identity
$\E[\A f]=0$ to get $b\bar v=\bar\lambda\mk$, and take expectations in
\eqref{eq:intensity} to get $\bar\lambda=\lambda_{0}+\lambda_{1}\bar v$;
solving the pair gives the display. The autocovariance follows from the
conditional mean by the tower property:
$\Cov(\Vj_{t},\Vj_{t+s})
=\E\big[(\Vj_t-\bar v)\,\E_t(\Vj_{t+s}-\bar v)\big]
=\Var(\Vj)e^{-\kappa_j s}$.
\end{proof}

\begin{remark}[Reading the theorem]
$b$ is the \emph{microscopic} decay: how fast one excitation wanes.
$\kappa_{j}=b(1-\eta)$ is the \emph{macroscopic} decay: how fast the whole
cascade relaxes, slowed by the factor $(1-\eta)$ because decaying excitation
keeps being replenished by offspring jumps. The data constrains
$\kappa_{j}$ (through every autocovariance and cross-covariance decay), while
$b$ appears separately only through the mean level $\bar v=\bar\lambda\mk/b$.
The Step-A transience requirement (``jump factor half-life under 2 days'')
properly applies to $b$; the persistence the diagnostics revealed lives in
$\kappa_{j}$. Both can now be honored at once.
\end{remark}

\begin{theorem}[Stationary variance]\label{thm:varvj}
$\displaystyle
S:=\Var(\Vj)=\frac{\bar\lambda\,\mkk}{2\kappa_{j}} .$
\end{theorem}

\begin{proof}
Take $f(v)=v^{2}$: $\A f(v)=-2bv^{2}
+(\lambda_{0}+\lambda_{1}v)\big(2v\mk+\mkk\big)$. Stationarity
($\E[\A f]=0$) gives
\[
2\big(b-\lambda_{1}\mk\big)\E[V^{2}]
=2\lambda_{0}\mk\,\bar v+\mkk\,\bar\lambda .
\]
From the proof of Theorem~\ref{thm:effective},
$\lambda_{0}\mk=\kappa_{j}\bar v$, so
$\E[V^{2}]=\bar v^{2}+\bar\lambda\mkk/(2\kappa_{j})$, and subtracting
$\bar v^{2}$ gives the claim.
\end{proof}

Note the same effective rate $\kappa_{j}$ divides the variance: stronger
self-excitation ($\eta\uparrow$) both slows the decay and fattens the
stationary spread --- one mechanism, two visible symptoms, which is why the
data can identify it.

% ===========================================================================
\chapter{The moment engine}
\label{ch:moments}

This chapter derives every population moment used downstream. Throughout,
$S=\Var(\Vj)$, $\kappa_{j}=b(1-\eta)$, $\bar\lambda$, $\bar v$ are as in
Chapter~\ref{ch:model}, and $\sigma_{c}^{2}=\sigma_{cv}^{2}\bar V^{c}/(2\kappa_{c})$
is the CIR stationary variance (unchanged from Step A).

\section{The damping trio}

Daily observables are \emph{time integrals}, which blur point-in-time
covariances. All blurring is captured by three scalar functions.

\begin{lemma}[Integration damping]\label{lem:damping}
Let $X$ be stationary with $\Cov(X_{t},X_{t+s})=\varsigma^{2}e^{-a|s|}$, and
let $I_{t}=\int_{t-1}^{t}X_{s}\dd s$. Define
\[
f(a)=\frac{2\,(a-1+e^{-a})}{a^{2}},\qquad
\hat g(a)=\frac{1-e^{-a}}{a},\qquad
g(a)=\hat g(a)^{2}=\frac{(1-e^{-a})^{2}}{a^{2}} .
\]
Then for integer $k\ge1$:
\[
\Var(I_{t})=\varsigma^{2}f(a),\qquad
\Cov(I_{t},I_{t-k})=\varsigma^{2}g(a)\,e^{-a(k-1)},\qquad
\Cov(X_{t},I_{t-k})=\varsigma^{2}\hat g(a)\,e^{-ak}\ (k\ge0).
\]
\end{lemma}

\begin{proof}
All three are Fubini integrals of $e^{-a|u-w|}$ over, respectively, the unit
square $\{u,w\in[0,1]\}$ (split along the diagonal, giving $f$), the
displaced square $\{u\in[k,k+1],w\in[0,1]\}$ (giving $g\,e^{-a(k-1)}$), and
the segment $\{w\in[k,k+1]\}$ with $u$ fixed at $0$ distance $k$ away (giving
$\hat g\,e^{-ak}$). The piecewise transition $f\ne g$ at the $k{=}0\to
k{=}1$ boundary is the structural discontinuity the spec calls the
``time-integration blur.''
\end{proof}

Numerically, $f,g,\hat g$ must be evaluated with \texttt{expm1} to avoid
catastrophic cancellation at small $a$ (Appendix~\ref{app:derivatives}
collects values, limits, and derivatives).

\section{The covariance kernel of jumps with the future state}

The single computation from which all cross moments flow:

\begin{lemma}[Feed-through kernel]\label{lem:kernel}
For $s\ge u$, with $M[\psi]$ the compensated $\psi$-jump integral,
\[
\Cov\!\big(\psi\text{-jump surprise in }\dd u,\ \Vj(s)\big)
=\bar\lambda\,m_{k\psi}\,e^{-\kappa_{j}(s-u)}\,\dd u ,
\]
and the covariance is zero for $s<u$. In integrated form, for the running
martingale $M_{s}=\int_{0}^{s}\!\int\psi\,\dd\tilde\mu$ started at $0$ from
stationarity,
$\ \Cov(\Vj(s),M_{s})
=\dfrac{\bar\lambda\,m_{k\psi}}{\kappa_{j}}\big(1-e^{-\kappa_{j}s}\big)$.
\end{lemma}

\begin{proof}
Apply Dynkin's formula on the fixed interval $[0,s]$ to the \emph{pair}
$(V,M)=(\Vj,M)$ with $f(v,m)=vm$. Between jumps, $\dd V=-bV\dd t$ and
$\dd M=-\lambda(t)m_{\psi}\dd t$ (the compensator drift); a jump with mark
$x$ moves $(v,m)\mapsto(v+k(x),\,m+\psi(x))$. Hence
\[
\A f(v,m) = -b\,vm \;-\;\lambda(v)\,m_{\psi}\,v
\;+\;\lambda(v)\big[v\,m_{\psi}+m\,\mk+m_{k\psi}\big]
= -b\,vm+\lambda(v)\big(m\,\mk+m_{k\psi}\big).
\]
Take expectations with $\phi(s):=\E[V(s)M(s)]$, use $\E[M]=0$ (so
$\E[\lambda M]=\lambda_{1}\E[VM]$) and $\E[\lambda]=\bar\lambda$:
\[
\phi(s)=\int_{0}^{s}\!\big[-\kappa_{j}\,\phi(u)
+\bar\lambda\,m_{k\psi}\big]\dd u
\;\Longrightarrow\;
\phi(s)=\frac{\bar\lambda m_{k\psi}}{\kappa_{j}}\big(1-e^{-\kappa_{j}s}\big),
\]
again a Volterra identity on a fixed interval. Differentiating the closed
form in $s$ and matching against
$\phi(s)=\int_{0}^{s}\text{(kernel)}(s,u)\dd u$ identifies the kernel as
$\bar\lambda m_{k\psi}e^{-\kappa_j(s-u)}$. For $s<u$
Lemma~\ref{lem:orthogonality} gives zero.
\end{proof}

Note carefully: the kernel decays at $\kappa_{j}$, \emph{not} $b$. A jump's
direct excitation decays at $b$, but its \emph{descendants} keep the
expectation elevated; the conditional-mean rate $\kappa_j$ governs
everything observable.

\section{Decomposing \texorpdfstring{$\JV$}{JV}}

The master decomposition, used everywhere below:
\begin{equation}\label{eq:jvdecomp}
\JV_{t}
=\underbrace{M_{t}}_{\text{jump surprise}}
+\;\mh\,\Lambda_{t},
\qquad
M_{t}:=\int_{t-1}^{t}\!\!\int h^{2}\dd\tilde\mu,
\quad
\Lambda_{t}:=\int_{t-1}^{t}\!\lambda(s)\dd s
=\lambda_{0}+\lambda_{1}\!\!\int_{t-1}^{t}\!\!\Vj(s)\dd s .
\end{equation}
$M_{t}$ is white (Lemma~\ref{lem:orthogonality}); all serial dependence of
$\JV$ rides on $\Lambda_{t}$, i.e.\ on the integrated $\Vj$ path --- which is
also what $\TPV$ observes. This one line already explains, qualitatively,
every diagnostic: $\JV$ autocorrelates because $\Lambda$ does; $\TPV$
predicts future $\JV$ because both load on $\Vj$; and jump surprises feed
future $\TPV$ through the kernel of Lemma~\ref{lem:kernel}.

\section{The complete moment set}

Write $\mathcal{E}_{k}:=e^{-\kappa_{j}(k-1)}$ and
$\mathcal{E}^{c}_{k}:=e^{-\kappa_{c}(k-1)}$ for $k\ge 1$.

\begin{theorem}[Population moments of the observables]\label{thm:moments}
Under Assumption~\ref{ass:model}, for lags $k\ge1$:
\begin{align}
&\textbf{Means:}&
\E[\TPV_{t}]&=\bar V^{c}+\bar v, &
\E[\JV_{t}]&=\bar\lambda\,\mh . \label{eq:means}\\[1ex]
&\textbf{TPV block:}&
\Var(\TPV_{t})&=\sigma_{c}^{2}f(\kappa_{c})+S\,f(\kappa_{j}),\hspace{-4em} && \label{eq:vartpv}\\
&&\Cov(\TPV_{t},\TPV_{t-k})&=\sigma_{c}^{2}g(\kappa_{c})\,\mathcal{E}^{c}_{k}
+S\,g(\kappa_{j})\,\mathcal{E}_{k}. \hspace{-4em}&&\label{eq:acovtpv}\\[1ex]
&\textbf{JV block:}&
\Var(\JV_{t})&=\bar\lambda\,\mhh
+\Big[\lambda_{1}^{2}\mh^{2}\,S+\lambda_{1}\mh\,\bar\lambda\,\mkh\Big]f(\kappa_{j}),
\hspace{-10em}&&\label{eq:varjv}\\
&&\Cov(\JV_{t},\JV_{t-k})&=
\underbrace{\lambda_{1}\mh\big[\lambda_{1}\mh\,S+\bar\lambda\,\mkh\big]}_{=:C_{J}}
\,g(\kappa_{j})\,\mathcal{E}_{k}. \hspace{-10em}&&\label{eq:acovjv}\\[1ex]
&\textbf{Cross block:}&
\Cov(\TPV_{t},\JV_{t})&=\bar\lambda\,\mkh\,\tfrac{1}{2}f(\kappa_{j})
+\lambda_{1}\mh\,S\,f(\kappa_{j}), \hspace{-10em}&&\label{eq:cross0}\\
&&\Cov(\TPV_{t},\JV_{t-k})&=
\underbrace{\big[\bar\lambda\,\mkh+\lambda_{1}\mh\,S\big]}_{=:C_{F}}
\,g(\kappa_{j})\,\mathcal{E}_{k}
\quad\text{(jumps lead),}\hspace{-12em}&&\label{eq:crossfwd}\\
&&\Cov(\TPV_{t},\JV_{t+k})&=
\underbrace{\lambda_{1}\mh\,S}_{=:C_{B}}
\,g(\kappa_{j})\,\mathcal{E}_{k}
\quad\text{(variance leads).}\hspace{-12em}&&\label{eq:crossbwd}
\end{align}
\end{theorem}

\begin{proof}
\emph{Means}: Lemma~\ref{lem:firstmoment} with $\psi=h^{2}$ and
Theorem~\ref{thm:effective}. \emph{TPV block}: $\Vc\perp\Vj$, each with an
exponential autocovariance ($\kappa_c$ from the CIR, $\kappa_j$ from
Theorem~\ref{thm:effective}); apply Lemma~\ref{lem:damping} factorwise.
\emph{JV block}: expand \eqref{eq:jvdecomp}.
$\Var(M_t)=\bar\lambda\mhh$ (Lemma~\ref{lem:isometry});
$\Var(\mh\Lambda_t)=\mh^2\lambda_1^2\,S f(\kappa_j)$
(Lemma~\ref{lem:damping} applied to $\Vj$);
$2\mh\Cov(M_t,\Lambda_t)
=2\mh\lambda_1\Cov\big(M_t,\int\Vj\big)
=2\mh\lambda_1\cdot\bar\lambda\mkh\cdot\tfrac12 f(\kappa_j)$, where the last
step integrates the kernel of Lemma~\ref{lem:kernel} ($\psi=h^2$,
$m_{kh^2}=\mkh$) over the triangle $\{u\le s\}$ of the unit square ---
exactly the $\tfrac12 f$ geometry of Lemma~\ref{lem:damping}. Summing gives
\eqref{eq:varjv}. For lags: $M$'s of disjoint days are mutually orthogonal
and orthogonal to the earlier path (Lemma~\ref{lem:orthogonality}), so only
$\Lambda$--$\Lambda$ and past-$M$--future-$\Lambda$ terms survive:
$\mh^2\lambda_1^2 S g(\kappa_j)\mathcal{E}_k$ and
$\mh\lambda_1\cdot\bar\lambda\mkh\,g(\kappa_j)\mathcal{E}_k$ (kernel
integrated over displaced squares); factoring yields \eqref{eq:acovjv}.
\emph{Cross block}: $\Cov(\int\Vj\text{-day }t,\ \JV_{t-k})$ has an
$M$-part (kernel over the displaced square:
$\bar\lambda\mkh\,g(\kappa_j)\mathcal{E}_k$) and a $\Lambda$-part
($\lambda_1\mh S g(\kappa_j)\mathcal{E}_k$), giving \eqref{eq:crossfwd};
at $k=0$ the $M$-part sees only the triangle ($\tfrac12 f$) while the
$\Lambda$-part sees the full square ($f$), giving \eqref{eq:cross0}; in the
backward direction the $M$-part dies (future surprises are orthogonal to the
past, Lemma~\ref{lem:orthogonality}) and only the $\Lambda$-part survives,
giving \eqref{eq:crossbwd}. $\Vc$ contributes nowhere in the JV/cross blocks
by independence.
\end{proof}

\begin{corollary}[Nesting]\label{cor:nesting}
At $\lambda_{1}=0$: $\kappa_{j}=b$, $\bar\lambda=\lambda_0$,
$C_{J}=C_{B}=0$, $C_{F}=\lambda_{0}\mkh$, and every formula collapses to the
Step-A model with $\mathcal{I}_{h^2}=\lambda_0\mh$,
$\mathcal{I}_{h^4}=\lambda_0\mhh$, $\mathcal{I}_{kh^2}=\lambda_0\mkh$.
\end{corollary}

\begin{corollary}[Structure of the extension]\label{cor:structure}
The three rejected diagnostics are repaired with signs and shapes forced by
the theory:
\begin{enumerate}
\item $\JV$ autocovariance is a single exponential in $\kappa_{j}$ with
positive weight $C_{J}$ --- matching the observed slow decay;
\item the cross-autocovariance is \emph{two-sided}, and its asymmetry is
exactly the co-jump coupling:
$C_{F}-C_{B}=\bar\lambda\,\mkh\ \ge 0$ --- matching the observed
``both sides positive, forward larger'';
\item all serial decay happens at the effective rate $\kappa_{j}$, which can
be slow (persistent cascades) even while $b$ is fast (transient individual
shocks) --- matching the week-long feed-through.
\end{enumerate}
Moreover the three coefficients obey the \textbf{internal consistency
identity}
\[
\boxed{\;C_{J}\cdot S \;=\; C_{B}\cdot C_{F}\;}
\]
(immediate from $C_J=\lambda_1\mh\,C_F$ and $C_B=\lambda_1\mh\,S$): an
over-identifying restriction imposed on \emph{empirical} quantities that the
data can reject. It is a free, sharp specification test and should be
reported with the calibration.
\end{corollary}

% ===========================================================================
\chapter{Identification and GMM design}
\label{ch:gmm}

\section{What is actually identified: composite parameters}

Inspect Theorem~\ref{thm:moments}: the raw quantities
$(\lambda_{0},\lambda_{1},\mh,\mhh,c)$ never appear alone. Every moment is a
function of the \emph{composites}
\[
A:=\bar\lambda\,\mh,\qquad
B:=\bar\lambda\,\mhh,\qquad
\psi:=\lambda_{1}\,\mh,\qquad
c,\qquad \kappa_{j},\qquad \kappa_{c},\ \sigma_{c}^{2},\ \bar V^{c},\ S,
\]
using the gauge to write $\mkh=c\,\mhh$, hence
$\bar\lambda\mkh=cB$, and (from Theorem~\ref{thm:effective})
\[
\bar v=\frac{\bar\lambda\,\mk}{b}=\frac{c\,A}{b},
\qquad
b=\kappa_{j}+\psi\,c,
\qquad
\eta=\frac{\psi c}{b}.
\]
In composite form the moment set reads:
\begin{align*}
\E[\JV]&=A, &
\Var(\JV)&=B+\big(\psi^{2}S+\psi\,cB\big)f(\kappa_{j}),\\
\E[\TPV]&=\bar V^{c}+\tfrac{cA}{\kappa_{j}+\psi c}, &
\gamma^{\JV}_{k}&=\psi\,(cB+\psi S)\;g(\kappa_{j})\,\mathcal{E}_{k},\\
\Var(\TPV)&=\sigma_{c}^{2}f(\kappa_{c})+S f(\kappa_{j}), &
\gamma^{\TPV}_{k}&=\sigma_{c}^{2}g(\kappa_{c})\mathcal{E}^{c}_{k}+S g(\kappa_{j})\mathcal{E}_{k},\\
\Cov(\TPV_t,\JV_t)&=cB\,\tfrac{f(\kappa_j)}{2}+\psi S f(\kappa_{j}), &
\text{cross}^{\pm}_{k}&=\big(cB+\psi S\big)\,\text{or}\,\big(\psi S\big)
\times g(\kappa_{j})\mathcal{E}_{k}.
\end{align*}

\begin{proposition}[Scale degeneracy and normalization]\label{prop:ident}
$(\lambda_{0},\lambda_{1},\mh,\mhh)$ enter only through $(A,B,\psi)$: the
split of jump activity between \emph{rate} and \emph{mark size} is not
identified from these moments (``twice as many jumps, half the energy each''
is observationally equivalent at this moment order). We therefore normalize
the mark scale, $\mh:=1$, after which
\[
\bar\lambda=A,\qquad \mhh=B/A,\qquad \lambda_{1}=\psi,
\qquad \lambda_{0}=\bar\lambda(1-\eta)
\]
are exactly recovered. This is a labeling convention, not a restriction:
every observable quantity and every downstream object (filter, VRP forecast)
depends on the composites only.
\end{proposition}

\begin{remark}[Gauge: strict vs.\ relaxed]\label{rem:gauge}
The gauge $k=c\,h^{2}$ applied to \emph{second} mark moments forces
$\mkk=c^{2}\mhh$, which by Theorem~\ref{thm:varvj} ties down
$S=c^{2}B/(2\kappa_{j})$ --- removing $S$ from the free set. We recommend
estimating with $S$ \textbf{free} (``relaxed gauge'': the gauge is used only
where the spec requires it, for the level recovery $\bar v=cA/b$), for two
reasons: it nests the Step-A treatment (where $\Var(V^j)$ was free), and it
converts the strict-gauge equation into a second testable over-identifying
restriction, $S\stackrel{?}{=}c^{2}B/(2\kappa_{j})$, to report alongside
$C_{J}S=C_{B}C_{F}$. If both restrictions survive, re-estimate under the
strict gauge for efficiency.
\end{remark}

\section{The estimable vector, bounds, and moment menu}

\begin{center}
\fbox{\parbox{0.92\textwidth}{
\textbf{Parameter vector (relaxed gauge, $\mh=1$):}
\[
\Theta'=\big[\;\kappa_{c},\ \sigma_{c}^{2},\ \bar V^{c},\
\kappa_{j},\ A,\ B,\ \psi,\ c,\ S\;\big]\in\R^{9}.
\]
\textbf{Derived (reported, not estimated):}\quad
$b=\kappa_{j}+\psi c$,\quad $\eta=\psi c/b$,\quad
$\bar\lambda=A$,\quad $\lambda_{0}=A(1-\eta)$,\quad
$\lambda_{1}=\psi$,\quad $\bar v=cA/b$,\quad
$\rho_{j}=-b$.
}}
\end{center}

\paragraph{Bounds.} All of
$\sigma_{c}^{2},\bar V^{c},A,B,\psi,c,S$ strictly positive; $\kappa_{c}$
in its persistence band as before. Two structural constraints deserve
emphasis:
\begin{itemize}
\item \textbf{Stationarity is free.} $\eta<1\iff\kappa_{j}>0$: by
parametrizing with $\kappa_{j}$ directly, subcriticality is enforced by a
simple positivity bound --- no nonlinear constraint needed. This is a real
design win of the composite parametrization.
\item \textbf{The old transience requirement moves to $b$.} ``Jump factor
half-life $<2$ days'' constrains $b=\kappa_{j}+\psi c\ge\ln 2/2$, a
\emph{linear} inequality in $(\kappa_{j},\psi,c)$-space (impose via
\texttt{scipy.optimize} linear constraint or check-and-report). $\kappa_{j}$
itself must be left \emph{free to be persistent} --- that freedom is the
entire point of the extension.
\end{itemize}

\paragraph{Moment menu.} The recommended GMM target set (all formulas from
Theorem~\ref{thm:moments}):
\begin{center}
\begin{tabular}{llr}
\toprule
block & moments & count\\
\midrule
means & $\E[\TPV],\ \E[\JV]$ & 2\\
variances & $\Var(\TPV),\ \Var(\JV)$ & 2\\
contemporaneous cross & $\Cov(\TPV_t,\JV_t)$ & 1\\
$\TPV$ autocovariance & $k=1,\dots,15$ & 15\\
$\JV$ autocovariance & $k=1,\dots,10$ & 10\\
forward cross (jumps lead) & $k=1,\dots,5$ & 5\\
backward cross (variance leads) & $k=1,\dots,5$ & 5\\
\midrule
total & & 40\\
\bottomrule
\end{tabular}
\end{center}
Nine parameters, forty moments: heavily over-identified, by design --- the
over-identification \emph{is} the specification test. Note the reversal from
Step A: the $\JV$ autocovariances, previously excluded because their
Jacobian rows were identically zero, are now the \emph{primary} carriers of
information about $\psi$ and $\kappa_{j}$; the backward cross moments are
the only block that separates $C_{B}=\psi S$ from $C_{F}=cB+\psi S$, i.e.\
they identify $c$.

\paragraph{Who identifies whom (the mental map).}
\begin{center}
\begin{tabular}{ll}
\toprule
parameter & primary identifying feature\\
\midrule
$\kappa_{j}$ & common decay slope of $\JV$-acov, cross-acov (log-linear in $k$)\\
$\kappa_{c},\sigma_{c}^{2}$ & slow residual component of $\TPV$-acov after the $\kappa_{j}$ part\\
$S$ & fast component weight in $\TPV$-acov ($S g(\kappa_{j})$)\\
$A$ & $\E[\JV]$\\
$B$ & $\Var(\JV)$ level (the $f$-terms are second order)\\
$\psi$ & $\JV$-acov weight $C_{J}$ and backward-cross weight $C_{B}=\psi S$\\
$c$ & cross asymmetry $C_{F}-C_{B}=cB$\\
$\bar V^{c}$ & $\E[\TPV]$ net of $\bar v=cA/b$\\
\bottomrule
\end{tabular}
\end{center}

\section{Weighting, conditioning, and the Jacobian}

Everything the Step-A code does survives, amplified:
\begin{itemize}
\item \textbf{Weighting matrix.} Forty moment contributions, many of them
products of heavy-tailed series, make the raw covariance of moment
conditions worse-conditioned than Step A's ($2.6\times10^{6}$ before
truncation). Ledoit--Wolf shrinkage remains mandatory; keep the
condition-number guardrail (warn above $100$, ridge backstop) and expect a
larger shrinkage intensity. Consider estimating $W$ once at the empirical
moments (as now) and, once stable, a second GMM step at $\hat\Theta'$.
\item \textbf{Analytic Jacobian.} Every moment is a product of (i) a
coefficient polynomial in $(\sigma_c^2,S,A,B,\psi,c)$ --- linear or
bilinear, trivial derivatives --- and (ii) a damping factor
$f,g,\hat g$ evaluated at $\kappa_{c}$ or $\kappa_{j}$ times an exponential
lag term. Chain rule through the composites; the only new scalar derivatives
needed are $f',g',\hat g'$ (Appendix~\ref{app:derivatives}) --- already in
the codebase for $f',g'$. The mean $\E[\TPV]$ additionally involves
$\partial_{(\kappa_j,\psi,c)}\,cA/(\kappa_{j}+\psi c)$: elementary
quotients.
\item \textbf{Multi-start.} The surface now has a genuine ridge:
$(\psi\uparrow, S\downarrow)$ trade off in $C_{B}$ while
$C_{F}$ compensates through $c$. Seed the multi-start grid from the
composite mental map above: $\hat\kappa_{j}$ from a log-linear fit to the
empirical $\JV$-acov slope, $\hat C_{B}/\hat S\to\psi$,
$(\hat C_{F}-\hat C_{B})/\hat B\to c$ --- i.e., \emph{method-of-moments
initialization from exactly the diagnostics that motivated the extension}.
\item \textbf{Uncertainty.} Parameters near bounds and ridge geometry make
point estimates insufficient: run the repository's block bootstrap over
$\Theta'$ and report intervals for the \emph{derived} quantities too
($\eta$, $b$, half-lives $\ln 2/\kappa_{j}$, $\ln 2/b$) --- the branching
ratio $\eta$ with its CI is the headline number of the whole exercise
(``how self-exciting is BTC jump risk?'').
\end{itemize}

% ===========================================================================
\chapter{Filtering: projecting the latent states}
\label{ch:filter}

\section{LMMSE without Gaussianity}

The Step-B objective is the best \emph{linear} projection of
$X_{t}=(\Vc(t),\Vj(t))'$ onto a window of observables
$Y_{t}=(\TPV_{t-j},\JV_{t-j})_{j=0}^{p-1}$:
\[
\hat X_{t}=\mu_{X}+\Sigma_{XY}\Sigma_{YY}^{-1}\,(Y_{t}-\mu_{Y}),
\qquad
\text{MSE}=\Sigma_{XX}-\Sigma_{XY}\Sigma_{YY}^{-1}\Sigma_{YX}.
\]
No Gaussianity is assumed or needed: among all linear functions of $Y_t$
this minimizes mean-squared error, period. What the model supplies is the
\emph{population second-moment structure}; better structure, better filter.
This is where the extension pays concretely: under Step A the model's
$\Sigma_{XY}$ said lagged $\JV$ was nearly worthless (loadings decaying at
$e^{-6.14k}$) --- under $\mathcal{M}'$ the loadings decay at
$e^{-\kappa_{j}k}$ with $\kappa_{j}$ estimated from data that clearly wants
persistence.

\section{The projection inputs, in closed form}

$\mu_{X}=(\bar V^{c},\ \bar v)'$; $\mu_{Y}$ stacks
\eqref{eq:means}. $\Sigma_{XX}=\operatorname{diag}(\sigma_{c}^{2},\,S)$
(the factors are independent). The cross block, from
Lemma~\ref{lem:damping} (third formula) and the kernel of
Lemma~\ref{lem:kernel} --- all with the \emph{uniform} lag law valid for
$k\ge0$:
\begin{equation}\label{eq:sigmaxy}
\boxed{
\begin{aligned}
\Cov\big(\Vc(t),\ \TPV_{t-k}\big)&=\sigma_{c}^{2}\ \hat g(\kappa_{c})\,e^{-\kappa_{c}k},
&\Cov\big(\Vc(t),\ \JV_{t-k}\big)&=0,\\
\Cov\big(\Vj(t),\ \TPV_{t-k}\big)&=S\ \hat g(\kappa_{j})\,e^{-\kappa_{j}k},
&\Cov\big(\Vj(t),\ \JV_{t-k}\big)&=C_{F}\ \hat g(\kappa_{j})\,e^{-\kappa_{j}k}.
\end{aligned}}
\end{equation}
(The $\Vj$--$\JV$ entry combines the jump-surprise kernel,
$\bar\lambda\mkh=cB$, with the intensity channel, $\psi S$; their sum is the
same $C_{F}$ as in the forward cross moment --- a pleasing and non-accidental
coincidence: the filter ``sees'' jump history through exactly the channel the
cross-covariance diagnostic measures.)

$\Sigma_{YY}$ stacks the full Theorem~\ref{thm:moments} moment set on the
chosen window.

\section{Model, hybrid, and the PSD guardrail}

Three ways to fill $(\Sigma_{XY},\Sigma_{YY})$:
\begin{enumerate}
\item \textbf{Pure model}: both from $\hat\Theta'$. Internally consistent;
optimal if $\mathcal{M}'$ is right; still misspecification-fragile.
\item \textbf{Hybrid}: $\Sigma_{XY}$ from the model, $\Sigma_{YY}$
empirical with Ledoit--Wolf shrinkage. Robust to residual $\Sigma_{YY}$
misspecification; the pragmatic default. \emph{Guardrail}: the stitched
joint matrix
$\bigl(\begin{smallmatrix}\Sigma_{XX}&\Sigma_{XY}\\ \Sigma_{YX}&\Sigma_{YY}\end{smallmatrix}\bigr)$
is not automatically PSD; check its smallest eigenvalue, and if negative,
shrink $\Sigma_{XY}\to\alpha\Sigma_{XY}$ with the largest
$\alpha\in(0,1]$ restoring PSD (report $\alpha$; $\alpha\ll1$ is itself a
misspecification alarm). Also verify the implied
$\text{MSE}$ matrix is PSD --- same check, same fix.
\item \textbf{Comparison as deliverable}: run both on simulated data from
$\mathcal{M}'$ (where the latent truth is known) and report state-recovery
RMSE. Under Step A's $\hat\Theta$ the pure-model filter is provably leaving
information on the table; \emph{quantifying} that gap is a genuinely
interesting result and directly showcases the misspecification-cost theme
of the write-up.
\end{enumerate}

\begin{impact}
Filter outputs $(\hat\Vc_{t},\hat\Vj_{t})$ are the state variables for
everything in Chapter~\ref{ch:vrp}. Their quality ceiling is set by the
covariances \eqref{eq:sigmaxy}; those are trustworthy only if the moment
formulas behind them survived (i) simulation verification
(Chapter~\ref{ch:impl}) and (ii) the over-identification tests of
Chapter~\ref{ch:gmm}. That chain is why I want the estimation nailed down before making any claim
about the VRP.
\end{impact}

% ===========================================================================
\chapter{The variance risk premium and regime classification}
\label{ch:vrp}

\section{Physical-measure forecasts in closed form}

Total quadratic variation over a horizon $\tau$ (in days) is
$\QV_{t,t+\tau}=\int_{t}^{t+\tau}(\Vc+\Vj)\dd s+\int\!\!\int h^{2}\dd\mu$.
Using the conditional means (Theorem~\ref{thm:effective} and its CIR
analogue) and Lemma~\ref{lem:firstmoment}, with
$\tilde\tau(a):=\big(1-e^{-a\tau}\big)/a$:
\begin{equation}\label{eq:pforecast}
\boxed{\;
\E^{\mathbb{P}}_{t}\big[\QV_{t,t+\tau}\big]
=\underbrace{\big(\bar V^{c}+\bar v+A\big)\tau}_{\text{unconditional level}}
+\big(\Vc_{t}-\bar V^{c}\big)\,\tilde\tau(\kappa_{c})
+\big(1+\psi\big)\big(\Vj_{t}-\bar v\big)\,\tilde\tau(\kappa_{j}) . \;}
\end{equation}
Two structural remarks. First, the jump-energy term enters through
$\E_{t}\int\lambda\,\mh$, and since $\lambda$ is affine in $\Vj$ it
\emph{adds} $\psi$ to the loading on the jump state: elevated jump variance
predicts future variance twice --- through its own decay \emph{and} through
the extra jumps it will cause. Second, the loading decays with
$\tilde\tau(\kappa_{j})$: for a 30-day horizon and a persistent
$\kappa_{j}$, jump-state shocks move even month-ahead expected variance
materially. Both effects are absent from the Step-A model. In practice,
plug the filtered $(\hat\Vc_{t},\hat\Vj_{t})$ into \eqref{eq:pforecast}.

\section{The premium is affine in the states}

Under a standard affine change of measure (intensity scaling
$\lambda^{\mathbb{Q}}=\ell\,\lambda$, mark tilting
$G^{\mathbb{Q}}$, and drift adjustments to $\Vc$), the
$\mathbb{Q}$-expectation of $\QV$ has the \emph{same functional form} as
\eqref{eq:pforecast} with risk-adjusted coefficients. Hence
\begin{equation}\label{eq:vrpaffine}
\mathrm{VRP}_{t}(\tau)
:=\E^{\mathbb{Q}}_{t}[\QV_{t,t+\tau}]-\E^{\mathbb{P}}_{t}[\QV_{t,t+\tau}]
=\alpha_{0}(\tau)+\alpha_{c}(\tau)\,\Vc_{t}+\alpha_{j}(\tau)\,\Vj_{t} :
\end{equation}
\emph{the variance risk premium is an affine function of the two latent
states.} This is the theoretical license for the entire classification
program: states of the VRP \emph{are} regions of the
$(\Vc,\Vj)$-plane, so classifying filtered states classifies the premium.

\paragraph{Getting the $\mathbb{Q}$-side without an option-pricing model.}
The clean two-step exploits \eqref{eq:vrpaffine}: obtain a model-free
$\mathbb{Q}$-expected variance from the options market (Deribit DVOL$^{2}$,
30-day), then regress it on the filtered states:
\[
\mathrm{DVOL}^{2}_{t}\cdot\tfrac{30}{365}
=\beta_{0}+\beta_{c}\hat\Vc_{t}+\beta_{j}\hat\Vj_{t}+\varepsilon_{t}.
\]
The theory says this linear regression is correctly specified (the
projection coincides with the structural relation up to filter error), so
$\hat\beta-$(the $\mathbb{P}$-coefficients from \eqref{eq:pforecast})
estimates the $\alpha(\tau)$'s directly. No parametric option model, no
risk-price calibration: honest, simple, and defensible in an interview.
Expected sign restrictions ($\beta_{j}$ exceeding the physical loading
$(1+\psi)\tilde\tau(\kappa_j)$, i.e.\ $\alpha_j>0$: jump risk carries a
premium) become testable hypotheses --- report them.

\section{Regimes and stability, made precise}

The affine structure tells you \emph{which} features are sufficient and
\emph{what stability means}:
\begin{itemize}
\item \textbf{Features.} The jump share $\hat\Vj_{t}/(\hat\Vc_{t}+\hat\Vj_{t})$
and the total level $\hat\Vc_{t}+\hat\Vj_{t}$ span the state plane; by
\eqref{eq:vrpaffine} any VRP regime is a half-plane-ish region in these
coordinates. Add the \emph{intensity innovation}
$\JV_{t}-\E[\JV_{t}\mid\text{history}]
=\JV_{t}-\big(A+\psi\,(\widehat{\Vj\text{-path}}-\bar v)\text{-term}\big)$:
under $\mathcal{M}'$ it is white, so its running size is a real-time
misspecification/stress alarm --- clustering that even the Hawkes term
cannot absorb.
\item \textbf{Classification.} Start with transparent thresholds on
(jump share, level); add a 2--3 state HMM on logs as a comparison; select
state count by held-out likelihood, not in-sample BIC alone, given how few
daily observations there are. Validate against the known event: 2025-10-10 must be flagged
without being told.
\item \textbf{Stability metrics with model meaning.} (i) Regime persistence
compares directly to the structural half-lives $\ln2/\kappa_{c}$,
$\ln2/\kappa_{j}$: a ``regime'' shorter than the states' own relaxation
times is noise, not regime. (ii) The branching ratio $\eta$ is a
\emph{global} fragility dial: as $\eta\to1$ the system approaches
criticality --- variance of $\Vj$ and cluster sizes diverge. Reporting
$\hat\eta$ with a bootstrap CI, and tracking it over subsamples
(2025 vs 2026), turns ``stability of the VRP'' from a vibe into a
parameter. (iii) Posterior entropy of the HMM gives a day-by-day
uncertainty measure; spikes flag transitions.
\end{itemize}

\begin{impact}
What this chapter makes possible, in order: the VRP time series
\eqref{eq:vrpaffine} on real data (needs the DVOL leg); regime-shaded charts
of VRP and jump share; the $\hat\eta$ criticality dial with CIs; an event
study around 2025-10-10 (VRP spike and decay at rate $\kappa_{j}$ ---
a directly testable model prediction: the post-event decay slope in the
data \emph{should} match the GMM-estimated $\kappa_{j}$; verifying it on
an event no moment was fitted to is the strongest check available here).
\emph{This is the prediction that failed when I ran it: observed persistence
was about 9.5 days against the 2.4 the fitted $\kappa_{j}$ implies.}
\end{impact}

% ===========================================================================
\chapter{Implementation: simulation, verification, and code changes}
\label{ch:impl}

\section{Exact simulation (no discretization)}

Two exact schemes; implement the first, keep the second as cross-check.

\paragraph{Ogata thinning.} Between jumps, $\lambda(t)$ decays
deterministically: from state $\Vj(\tau_{i})$,
$\lambda(t)=\lambda_{0}+\lambda_{1}\Vj(\tau_{i})e^{-b(t-\tau_{i})}$, which is
bounded by its left value. Iterate: (1) set
$\bar\lambda_{\text{loc}}=\lambda(\tau_{i}^{+})$; (2) propose the next event
after an $\mathrm{Exp}(\bar\lambda_{\text{loc}})$ wait; (3) accept with
probability $\lambda(\text{proposed time})/\bar\lambda_{\text{loc}}$
(the decayed, hence smaller, true rate); if rejected, update the bound and
repeat; (4) on acceptance draw $x\sim G$, add $k(x)$ to $\Vj$, record
$h^{2}(x)$. This is exact for monotone-between-jumps intensities. Simulate
$\Vc$ with the standard exact CIR transition (noncentral $\chi^{2}$) on a
fine grid; form daily $\TPV$, $\JV$ by integrating/aggregating.

\paragraph{Cluster construction.} Draw immigrants as Poisson($\lambda_{0}$)
on $[0,T]$; recursively give each event $\mathrm{Poisson}(\lambda_{1}k(x)/b)$
children with ages $\sim\mathrm{Exp}(b)$; stop when a generation is empty
(a.s.\ finite by $\eta<1$). Statistically identical output, independent code
path --- ideal for catching bugs in either.

\section{The verification protocol}

Before the first GMM run on real data:
\begin{enumerate}
\item Fix a test parameter point with meaningful self-excitation
(e.g.\ $\eta\approx0.6$, $\kappa_{j}$ half-life $\approx4$ d, $b$ half-life
$\approx0.3$ d).
\item Simulate $\ge2\times10^{5}$ days; compute all 40 empirical moments.
\item Compare against Theorem~\ref{thm:moments} with batch-means Monte Carlo
error bands; every formula must land inside its band. Repeat at
$\lambda_{1}=0$ and confirm collapse to the Step-A formulas
(Corollary~\ref{cor:nesting}).
\item Verify the analytic Jacobian against finite differences at random
interior points (max relative error $<10^{-6}$ at sensible step sizes).
\item Only then: run GMM on simulated data and confirm parameter recovery
(estimates within bootstrap bands of truth). This end-to-end rehearsal is
also Figure~1 of the write-up's appendix: ``the estimator works when the
model is true.''
\end{enumerate}

\section{Concrete changes to the codebase}

\begin{center}
\small
\begin{tabular}{lll}
\toprule
component & Step A & Step A$'$\\
\midrule
data engineering & 5-min bars & 1-min bars (signature-justified)\\
jump-robust estimator & tripower, $\min(\cdot,\RV)$ cap & same, add MedRV cross-check\\
\texttt{theoretical\_moments} & 20 moments, 8 params & 40 moments, 9 params ($\Theta'$)\\
damping functions & $f,g$ & $f,g,\hat g$ and derivatives\\
weighting & Ledoit--Wolf + ridge & shrink to \emph{diagonal}, not identity\\
bounds & $\kappa_c$ band, $\rho_j$ band & $\kappa_c$ band, positivity, $b\ge\ln2/2$\\
new module & --- & simulator + moment verifier\\
new outputs & 3 fit plots & + $\JV$-acov and two-sided cross fits,\\
& & tests $C_JS=C_BC_F$, strict gauge\\
filter (Step B) & --- & $\Sigma_{XY}$ \eqref{eq:sigmaxy}; hybrid $\Sigma_{YY}$; PSD guard\\
\bottomrule
\end{tabular}
\end{center}

Table~\ref{tab:dictionary} is the parameter dictionary. I keep it in the
repository docs and in the docstrings, with each block citing the equation
number it comes from.

\begin{table}[h]
\centering
\caption{Dictionary: Step-A parameters vs.\ Step-A$'$ composites.}
\label{tab:dictionary}
\begin{tabular}{lll}
\toprule
Step A & Step A$'$ & note\\
\midrule
$\mathcal{I}_{h^2}$ & $A=\bar\lambda\mh$ & $=\E[\JV]$\\
$\mathcal{I}_{h^4}$ & $B=\bar\lambda\mhh$ & $=\Var(M)$ level\\
$\mathcal{I}_{kh^2}$ & $cB=\bar\lambda\mkh$ & cross asymmetry $C_F-C_B$\\
$\Var(V^j)$ & $S$ & free under relaxed gauge\\
$-\rho_{j}$ & $b=\kappa_{j}+\psi c$ & microscopic decay (transient)\\
--- & $\kappa_{j}=b(1-\eta)$ & macroscopic decay (persistent)\\
--- & $\psi=\lambda_{1}\mh$, $\eta=\psi c/b$ & self-excitation strength\\
$\kappa_{c},\ \sigma_{c}^{2}$, level & unchanged & CIR block untouched\\
\bottomrule
\end{tabular}
\end{table}

% ===========================================================================
\appendix

\chapter{Damping functions and derivatives (for the analytic Jacobian)}
\label{app:derivatives}

With $E:=e^{-a}$ (evaluate $1-E$ as \texttt{-expm1(-a)} and
$a-1+E$ as \texttt{a + expm1(-a)}):
\begin{align*}
f(a)&=\frac{2(a-1+E)}{a^{2}}, &
f'(a)&=\frac{2(1-E)}{a^{2}}-\frac{4(a-1+E)}{a^{3}},\\
\hat g(a)&=\frac{1-E}{a}, &
\hat g'(a)&=\frac{E}{a}-\frac{1-E}{a^{2}},\\
g(a)&=\hat g(a)^{2}, &
g'(a)&=2\,\hat g(a)\,\hat g'(a) .
\end{align*}
Small-$a$ limits (useful as unit tests): $f(a)\to1-\tfrac{a}{3}$,
$\hat g(a)\to1-\tfrac{a}{2}$, $g(a)\to1-a$. Large-$a$:
$f\sim2/a$, $\hat g\sim1/a$, $g\sim1/a^{2}$. Lagged terms: for
$m(a)=g(a)e^{-a(k-1)}$, $m'(a)=\big[g'(a)-(k-1)g(a)\big]e^{-a(k-1)}$, and
analogously with $\hat g$ and lag $k$ ($\hat m'=[\hat g'-k\hat g]e^{-ak}$).
Derived-parameter chain rules: with $b=\kappa_{j}+\psi c$,
$\partial b/\partial\kappa_{j}=1$, $\partial b/\partial\psi=c$,
$\partial b/\partial c=\psi$; with $\bar v=cA/b$,
$\partial\bar v = \big(A/b,\ cA\,\partial(1/b)\big)$ componentwise.

\chapter{Useful integrals}
\label{app:integrals}
For $a>0$, $k\ge1$ integer:
\[
\int_{0}^{1}\!\!\int_{0}^{1}e^{-a|u-w|}\dd u\dd w=f(a),\qquad
\int_{k}^{k+1}\!\!\int_{0}^{1}e^{-a(s-u)}\dd u\dd s=g(a)e^{-a(k-1)},
\]
\[
\int_{0}^{1}\!\!\int_{0}^{s}e^{-a(s-u)}\dd u\dd s=\tfrac{1}{2}f(a),\qquad
\int_{k}^{k+1}e^{-as}\dd s=\hat g(a)e^{-ak},\qquad
\int_{0}^{\tau}e^{-as}\dd s=\tilde\tau(a)=\frac{1-e^{-a\tau}}{a}.
\]
The triangle integral (third) is exactly half the square integral (first):
the geometric reason the contemporaneous cross moment carries $f/2$ where
the variance carries $f$.

\chapter{Notation glossary}
\label{app:notation}
\begin{center}
\begin{tabular}{ll}
\toprule
$\mu,\ \tilde\mu,\ \nu$ & jump counting measure, compensated measure, compensator\\
$G,\ m_{\phi}$ & mark distribution (probability); $m_{\phi}=\int\phi\,\dd G$\\
$\lambda_{0},\lambda_{1},\lambda(t)$ & baseline, excitation slope, intensity $\lambda_0+\lambda_1\Vj(t^-)$\\
$b=-\rho_{j}$ & microscopic (per-jump) decay rate of $\Vj$\\
$\eta=\lambda_{1}\mk/b$ & branching ratio (children per jump); $\eta<1$\\
$\kappa_{j}=b(1-\eta)$ & effective (observable) decay rate of $\Vj$\\
$\bar\lambda,\ \bar v$ & stationary mean intensity, stationary mean of $\Vj$\\
$S$ & $\Var(\Vj)=\bar\lambda\mkk/(2\kappa_{j})$\\
$A,B,\psi,c$ & composites $\bar\lambda\mh$, $\bar\lambda\mhh$, $\lambda_{1}\mh$, gauge constant\\
$C_{F},C_{B},C_{J}$ & cross-forward, cross-backward, $\JV$-acov coefficients\\
$f,g,\hat g$ & damping trio (Lemma~\ref{lem:damping}); $g=\hat g^{2}$\\
$M_{t},\Lambda_{t}$ & day-$t$ jump surprise, integrated intensity \eqref{eq:jvdecomp}\\
$\mathcal{E}_{k}$ & $e^{-\kappa_{j}(k-1)}$ (and $\mathcal{E}^{c}_{k}$ with $\kappa_{c}$)\\
\bottomrule
\end{tabular}
\end{center}

\chapter{Implementation checklist}
\label{app:checklist}
\begin{enumerate}
\item Switch data engineering to 1-minute bars; add MedRV alongside tripower;
keep the $\TPV\le\RV$ truncation and its warning.
\item Implement $f,g,\hat g$ + derivatives with \texttt{expm1}; unit-test
against the small-/large-$a$ limits of Appendix~\ref{app:derivatives}.
\item Implement \texttt{theoretical\_moments\_ext($\Theta'$)} returning the
40-vector of Theorem~\ref{thm:moments} and its analytic $40\times9$
Jacobian; unit-test the $\lambda_{1}=0$ collapse against the existing
Step-A function.
\item Build the simulator (thinning; cluster construction as cross-check);
run the five-step verification protocol of Chapter~\ref{ch:impl}. Do not
proceed past a failing band.
\item Extend \texttt{empirical\_moments} to the 40-moment menu; reuse
Ledoit--Wolf weighting with condition-number guardrails.
\item Estimate; report $\Theta'$, derived $(b,\eta,\lambda_{0},\bar v)$ with
block-bootstrap CIs; report the consistency tests $C_{J}S=C_{B}C_{F}$ and
$S=c^{2}B/(2\kappa_{j})$; regenerate all three diagnostic figures with the
new model-implied curves overlaid.
\item Step B: assemble $\Sigma_{XY}$ from \eqref{eq:sigmaxy}, hybrid
$\Sigma_{YY}$, PSD guardrail; simulation study comparing pure-model vs
hybrid filters under both $\hat\Theta$ (old) and $\hat\Theta'$ (new).
\item VRP: source DVOL; regress on filtered states per
Chapter~\ref{ch:vrp}; build the VRP series, regimes, $\hat\eta$ dial, and
the 2025-10-10 event study (test: post-event decay slope $=\hat\kappa_{j}$).
\end{enumerate}

\end{document}
